\section{Fixed physical ports and bounded local registers}

The intersite gates of~\cite[Supplement, pp.~7--8]{Piroli2021LOCC}
act on fixed-dimensional physical qudits. Local ancillas are separate:
free onsite operations may mix the physical qudit with its local ancillas,
but a counted intersite gate does not act on an arbitrarily enlarged pair
of sites. This section derives a resource bound for encoded matching
channels in that convention. It does not identify the unrestricted
local-channel model with the source's phase relation; see
\cite{gap:psc21_physical_port_simulation_scope}.

\begin{definition}[Physical-port channel protocols]
    \label{def:qc_physical_port_protocol}
    \lean{QuantumCircuit.PhysicalPortLayout,
        QuantumCircuit.IsPhysicalPortUnitary,
        QuantumCircuit.IsPhysicalPortUnitary.mem_unitary,
        QuantumCircuit.IsPhysicalPortUnitary.conjTranspose,
        QuantumCircuit.IsPhysicalPortUnitary.one,
        QuantumCircuit.IsPhysicalPortUnitary.onsite_swap,
        QuantumCircuit.IsPhysicalPortProtocol,
        QuantumCircuit.IsPhysicalPortProtocol.isKrausCPTP,
        QuantumCircuit.IsPhysicalPortProtocol.one}
    \leanok
    Every wire carries a $d$-level qudit. A layout assigns its wires to spatial
    sites and selects one physical port at each site. Free operations are
    unitaries supported inside one spatial site, wire permutations preserving
    each spatial site, and normalized onsite Kraus channels. A counted layer
    is an existing nearest-neighbor unitary layer placed on the physical ports,
    with identity on every memory wire. Products add physical depth bounds;
    a larger upper bound remains valid. Reversing a unitary circuit preserves
    its depth. Every resulting channel protocol is CPTP on all input operators.
\end{definition}

\begin{lemma}[Register codes and completed decoding]
    \label{lem:qc_register_code_completion}
    \lean{QuantumCircuit.registerEncoding,
        QuantumCircuit.registerEncoding_isIsometry,
        QuantumCircuit.registerEncoding_isKrausCPTP,
        QuantumCircuit.registerDecoding,
        QuantumCircuit.registerDecoding_isKrausCPTP,
        QuantumCircuit.registerDecoding_encoding,
        QuantumCircuit.registerDecoding_comp_encoding,
        QuantumCircuit.registerDecoding_encoding_reference,
        QuantumCircuit.registerEncodedMap,
        QuantumCircuit.registerEncodedMap_isKrausCPTP,
        QuantumCircuit.registerEncodedMap_encoding,
        QuantumCircuit.le_registerDimension,
        QuantumCircuit.registerBasisEmbedding,
        QuantumCircuit.pairRegisterBasisEmbedding,
        QuantumCircuit.pairRegisterBasisEmbedding_apply,
        QuantumCircuit.registerEncoding_apply,
        QuantumCircuit.registerEncoding_pair}
    \leanok
    An injection $e$ of local basis labels into register words defines an
    isometry $E$. The encoding is $X\mapsto EXE^\dagger$. For a fixed density
    matrix $\rho$, the completed decoder is
    \begin{align}
        D(Y) & =E^\dagger YE+
        \operatorname{tr}[(I-EE^\dagger)Y(I-EE^\dagger)]\rho.
    \end{align}
    It is CPTP and satisfies $D(EXE^\dagger)=X$ for every operator $X$.
    The same identity holds after tensoring with the identity on any finite
    reference. Encoding the two endpoints separately gives the tensor product
    of their local encoding matrices. If $m\le B$ and $d\ge2$, then
    $k=\lceil\log_d B\rceil$ qudits suffice. A channel $\Phi$ extends to the
    full input register by $\widehat\Phi=\mathcal E_{\rm out}\Phi D_{\rm in}$;
    this extension is CPTP and agrees with $\Phi$ on encoded inputs.
\end{lemma}
\begin{proof}
    \leanok
    The inclusion has orthonormal columns. Complete the compression outside its
    range using the existing support-completion channel; the complementary
    term vanishes on $EXE^\dagger$. Matrix multiplication gives the exact left
    inverse, and tensoring preserves composition. The paired basis inclusion
    factors entrywise into two local inclusions. The ceiling-logarithm bound
    is $B\le d^k$.
\end{proof}

\begin{lemma}[Parallel physical SWAP layers]
    \label{lem:qc_port_matching_swaps}
    \uses{def:qc_physical_port_protocol}
    \lean{QuantumCircuit.permOpHom,
        QuantumCircuit.permOp_one,
        QuantumCircuit.permOp_mul,
        QuantumCircuit.permOp_conjTranspose,
        QuantumCircuit.PortMatching.perm,
        QuantumCircuit.PortMatching.perm_apply_of_notMem,
        QuantumCircuit.PortMatching.perm_apply_of_mem,
        QuantumCircuit.PortMatching.perm_apply_left,
        QuantumCircuit.PortMatching.perm_apply_right,
        QuantumCircuit.PortMatching.perm_involutive,
        QuantumCircuit.PortMatching.perm_symm,
        QuantumCircuit.PortMatching.layer,
        QuantumCircuit.PortMatching.layer_op,
        QuantumCircuit.PortMatching.layer_op_one_site}
    \leanok
    A matching of ring bonds defines an involutive site permutation: each
    selected bond exchanges its endpoints and every other site is fixed.
    The actual layer containing these SWAP gates has exactly the associated
    permutation matrix. Permutation matrices preserve products and inverses.
    The self-loop of a one-site ring acts as the identity.
\end{lemma}
\begin{proof}
    \leanok
    Swaps on disjoint bonds commute and fix each other's endpoints. Their
    finite product therefore has the stated action. Applying the permutation
    matrix homomorphism identifies that product with the layer operator.
\end{proof}

\begin{theorem}[Selected-memory routing at unchanged physical depth]
    \label{thm:qc_port_register_routing}
    \uses{def:qc_physical_port_protocol, lem:qc_port_matching_swaps}
    \lean{QuantumCircuit.PhysicalPortLayout.selectionPermutation,
        QuantumCircuit.PhysicalPortLayout.selectionPermutation_site,
        QuantumCircuit.PhysicalPortLayout.selectionPermutation_port,
        QuantumCircuit.PhysicalPortLayout.selectionPermutation_symm,
        QuantumCircuit.IsPhysicalPortUnitary.selected_layer,
        QuantumCircuit.embedOp_permOp_viaEmbedding,
        QuantumCircuit.IsPhysicalPortUnitary.selected_permutation_layer,
        QuantumCircuit.registerPermutation,
        QuantumCircuit.registerPermutation_apply,
        QuantumCircuit.registerPermutation_apply_of_notMem,
        QuantumCircuit.registerPermutation_involutive,
        QuantumCircuit.IsPhysicalPortUnitary.register_matching}
    \leanok
    Select one memory wire at each site. Swap it with the site's physical port,
    apply a physical layer, and undo the onsite swaps. This acts exactly as the
    layer on the selected wires and as identity on all other wires, at physical
    depth one. Repeating this for $k$ disjoint register digits implements the
    same matching permutation on each digit at depth at most $k$, independently
    of the number of selected bonds. All unselected wires are restored on the
    whole operator space, without assuming they were initialized or uncorrelated.
\end{theorem}
\begin{proof}
    \leanok
    The selection permutation preserves each spatial site, hence is free.
    Conjugating the port layer by it changes only the wire placement. Extend
    each digit's site permutation by identity outside that digit's selected
    wires. Disjoint ranges ensure the finite product has the desired action
    on each digit; every factor has a one-layer implementation.
\end{proof}

\begin{lemma}[Data and scratch placement]
    \label{lem:qc_port_data_scratch}
    \uses{thm:qc_port_register_routing}
    \lean{QuantumCircuit.PortRegisters.layout,
        QuantumCircuit.PortRegisters.wire,
        QuantumCircuit.PortRegisters.site_wire,
        QuantumCircuit.PortRegisters.data,
        QuantumCircuit.PortRegisters.scratch,
        QuantumCircuit.PortRegisters.site_data,
        QuantumCircuit.PortRegisters.site_scratch,
        QuantumCircuit.PortRegisters.data_ne_scratch,
        QuantumCircuit.PortRegisters.selection,
        QuantumCircuit.PortRegisters.selection_apply,
        QuantumCircuit.PortRegisters.site_selection,
        QuantumCircuit.PortRegisters.selection_ranges_disjoint,
        QuantumCircuit.PortRegisters.pairData,
        QuantumCircuit.PortRegisters.pairData_left,
        QuantumCircuit.PortRegisters.pairData_right,
        QuantumCircuit.PortMatching.right_notMem,
        QuantumCircuit.PortRegisters.routingPermutation,
        QuantumCircuit.PortRegisters.data_notMem_selection,
        QuantumCircuit.PortRegisters.routingPermutation_data_left,
        QuantumCircuit.PortRegisters.routingPermutation_data_of_notMem,
        QuantumCircuit.PortRegisters.routingPermutation_pairData_site,
        QuantumCircuit.PortRegisters.routingPermutation_isPhysicalPortUnitary,
        QuantumCircuit.PortRegisters.pairData_isPhysicalPortProtocol}
    \leanok
    Use $1+2k$ wires per spatial site: one physical port, $k$ data qudits and
    $k$ scratch qudits. On a matching, select a data digit at each left endpoint
    and a scratch digit at each right endpoint. The common routing moves the
    left data register into the right scratch register and leaves the right
    data register fixed. Thus both data registers of each nontrivial selected
    bond lie at its right spatial site after routing. Reversing the routing
    returns their outputs to their original sites and restores all other wires.
\end{lemma}
\begin{proof}
    \leanok
    A right endpoint of a nontrivial matching bond cannot be another left
    endpoint. Data and scratch slots are disjoint, and different digit selections
    have disjoint ranges. The action formulas of
    Theorem~\ref{thm:qc_port_register_routing} give the claimed placement.
\end{proof}

\begin{theorem}[One shared routing for a channel family]
    \label{thm:qc_port_shared_channel_routing}
    \uses{def:qc_physical_port_protocol, thm:qc_port_register_routing}
    \lean{QuantumCircuit.sum_conjTranspose_embedOp_mul,
        QuantumCircuit.routingMatrixEquiv,
        QuantumCircuit.routingMatrixEquiv_apply,
        QuantumCircuit.routingMatrixEquiv_symm_apply,
        QuantumCircuit.routingChannelConj,
        QuantumCircuit.routingChannelConj_apply,
        QuantumCircuit.routingChannelConj_embeddedKraus,
        QuantumCircuit.IsPhysicalPortProtocol.of_routed_map,
        QuantumCircuit.IsPhysicalPortProtocol.of_register_routing,
        QuantumCircuit.IsPhysicalPortProtocol.of_routed_family}
    \leanok
    Suppose an actual wire routing of physical depth $T$ moves every channel
    in a commuting family into a single spatial site, possibly a different
    site for each member. Execute all those conjugated channels locally,
    then reverse the routing. The result is exactly the original family
    product, with physical depth at most $2T$. The number of family members
    does not multiply this depth bound. The identities hold for every input
    operator, including arbitrary correlations with unused registers.
\end{theorem}
\begin{proof}
    \leanok
    Wire reindexing is a linear equivalence on matrices. Its conjugation on
    linear maps is multiplicative, so it preserves the commuting family
    product. Placing Kraus operators preserves their completeness relation.
    Each routed map is a zero-depth onsite channel; only the one forward and
    one backward routing contribute to the depth.
\end{proof}

\begin{theorem}[Uniform bounded-register matching simulation]
    \label{thm:qc_bounded_physical_port_simulation}
    \uses{lem:qc_register_code_completion, lem:qc_port_data_scratch,
        thm:qc_port_shared_channel_routing}
    \lean{QuantumCircuit.PortRegisters.pairChannel,
        QuantumCircuit.PortRegisters.pairChannel_commute,
        QuantumCircuit.PortRegisters.matchingChannel,
        QuantumCircuit.PortRegisters.matchingChannel_isPhysicalPortProtocol,
        QuantumCircuit.boundedPairRegisterCode,
        QuantumCircuit.exists_bounded_matching_simulation,
        QuantumCircuit.IsPhysicalPortProtocol.list_prod_uniform}
    \leanok
    Fix $d\ge2$, $B$, and a local input dimension $m\le B$, and choose a
    fallback density matrix on a pair of local inputs. Set
    $k=\lceil\log_d B\rceil$, independently of the chain length.
    For every ring of $N\ge2$ sites and every matching of CPTP pair channels,
    there are normalized Kraus maps on the two $k$-qudit data registers that
    agree with the prescribed code extensions on the entire register space,
    and intertwine exactly with the original pair channels on encoded inputs.
    Their actual placed matching channel has a physical-port protocol of depth
    at most $2k$. A product of $L$ such certified layers has depth at most $2kL$.
    For $B=1$ the register width and this bound are zero.
\end{theorem}
\begin{proof}
    \leanok
    Encode the two endpoints independently, then complete each pair channel
    outside its joint code. The resulting map is CPTP, so choose its normalized
    finite Kraus representation. The common data-to-scratch routing has depth
    $k$ and puts every pair channel onsite. Apply
    Theorem~\ref{thm:qc_port_shared_channel_routing} once to the complete
    matching. Protocol composition gives the final linear-in-$L$ bound.
    The off-code joint decoder need not factor onsite; it is only executed
    after colocation. This theorem is not an assertion about protocols with
    an unrestricted intermediate dimension depending on $N$.
\end{proof}

\begin{lemma}[Product onsite coding on arbitrary chain inputs]
    \label{lem:qc_product_register_encoding}
    \uses{lem:qc_register_code_completion}
    \lean{QuantumCircuit.OnsiteChannel.ofSiteMaps,
        QuantumCircuit.OnsiteChannel.ofSiteMaps_siteMap,
        QuantumCircuit.OnsiteChannel.comp,
        QuantumCircuit.OnsiteChannel.comp_map,
        QuantumCircuit.OnsiteChannel.map_eq_id_of_siteMap_eq_id,
        QuantumCircuit.OnsiteChannel.encodeRegisters,
        QuantumCircuit.OnsiteChannel.encodeRegisters_siteMap,
        QuantumCircuit.OnsiteChannel.encodeRegisters_map,
        QuantumCircuit.OnsiteChannel.decodeRegisters,
        QuantumCircuit.OnsiteChannel.decodeRegisters_encodeRegisters,
        QuantumCircuit.OnsiteChannel.decodeRegisters_encodeRegisters_reference,
        QuantumCircuit.OnsiteChannel.encoded,
        QuantumCircuit.OnsiteChannel.encoded_map_encodeRegisters}
    \leanok
    Finite CPTP maps at individual sites assemble into an onsite channel,
    and assembly preserves composition. Product isometric encodings have
    product completed local decoders whose composition is the identity on
    all chain operators, also with any finite reference. Extending an existing
    onsite channel between site-consistent input and output codes intertwines
    with its original chain map, including entangled chain inputs.
\end{lemma}
\begin{proof}
    \leanok
    Choose normalized local Kraus families. On product operators, all formulas
    follow site by site; these products span the entire chain matrix algebra.
    Use the exact local decoding identity before tensoring with a reference.
\end{proof}

\begin{lemma}[Representation-independent placement of register channels]
    \label{lem:qc_register_channel_lift}
    \uses{lem:qc_placed_operators}
    \lean{QuantumCircuit.registerConfigurationSplit,
        QuantumCircuit.registerConfigurationSplit_apply,
        QuantumCircuit.registerConfigurationSplit_symm_comp,
        QuantumCircuit.registerMatrixSplit,
        QuantumCircuit.registerMatrixSplit_conjTranspose,
        QuantumCircuit.registerMatrixSplit_embedOp,
        QuantumCircuit.registerChannelLift,
        QuantumCircuit.registerMatrixSplit_channelLift,
        QuantumCircuit.registerChannelLift_id,
        QuantumCircuit.registerChannelLift_comp,
        QuantumCircuit.registerChannelLift_kraus,
        QuantumCircuit.rectangularKrausMap_embedOp_congr}
    \leanok
    Splitting a configuration into a selected register and its complementary
    wires identifies a placed operator with $A\otimes I$. Under the same
    reindexing, placing a Kraus family gives precisely its local channel
    tensored with the identity on the complement. Consequently the placed
    channel is independent of its chosen Kraus representation, and the
    placement identity holds on arbitrary correlated inputs.
\end{lemma}
\begin{proof}
    \leanok
    Use the explicit configuration equivalence and the matrix reindexing
    algebra equivalence. Check tensoring on matrix units, then use linear
    extensionality. Applying the split to each placed Kraus operator proves
    the full channel identity.
\end{proof}

\begin{theorem}[Intertwining of placed encoded channels]
    \label{thm:qc_encoded_channel_placement}
    \uses{lem:qc_product_register_encoding, lem:qc_register_channel_lift}
    \lean{QuantumCircuit.reindex_rectKronecker_registerSplit,
        QuantumCircuit.registerMatrixSplit_encodeRegisters,
        QuantumCircuit.registerChannelLift_encodeRegisters}
    \leanok
    Let $E_S$ be the product code of a selected register and $E$ the product
    code of the complete chain, possibly with different local input and output
    dimensions. If $\widehat\Phi E_S=E_S\Phi$ as linear maps on operators,
    then the placed channels satisfy
    $\operatorname{Lift}(\widehat\Phi)E=E\operatorname{Lift}(\Phi)$
    on every chain operator. The remaining sites may carry arbitrary
    correlations with the selected register.
\end{theorem}
\begin{proof}
    \leanok
    The rectangular global encoding matrix splits as the tensor product of
    its selected and complementary encoding matrices. Check the resulting
    tensor intertwining on matrix units, then use linear extensionality and
    the representation-independent placement identity.
\end{proof}
